% s-185-p185-m002; original PDF 185, 186
\Aphorism{s-185-p185-m002}{275}
An Menschen, mit denen wir einmal flüchtig verkehrt haben, bewirkt die Erinnerung eine seltsame perspectivische Verzerrung. Einiges von ihnen behalten wir, einen Gesichtszug, eine Gebärde, ihren Tonfall oder Dialect oder eine ihrer Redensarten, während das Übrige durch das Sieb fällt; richten wir nun das Gedächtniss auf das hinterlassene Bild, so erscheint das Eine monströs nah und deutlich, der ganze Mensch aber fern, unglaubwürdig, traumhaft, und die Gesammtwirkung ist ein wenig lächerlich. Dieses Eine, die mnemonische Handhabe, an der wir den ganzen Menschen festhalten, braucht nicht einmal ein Theil von ihm selbst zu sein; oft bewahren wir ihn auf in Association mit einem gleichzeitigen Erlebniss und färben ihn mit dessen Gefühlsfarbe, die Lust oder Unlust bedeuten kann. Manchem Unschuldigen sind wir in der Erinnerung gram, nur weil wir zufällig Zahnweh hatten, als er uns in den Weg trat.
\par

